\section{The CNN: FoGNet Architecture}
\label{sec:cnn}

The shipped detector is a \emph{glass box}: one hand-chosen ratio, the freeze index,
compared against one hand-chosen trip line. That transparency is a feature, not a
bug --- a clinician can read off exactly why a window was flagged. But it also fixes
the rules in advance. The $3$--$8$\,\unit{\hertz} freeze band, the $0.5$--$3$\,\unit{\hertz}
locomotor band and the simple ``ratio of powers'' shape were all decided by us, not
discovered from data. \textbf{FoGNet} asks the opposite question: \emph{given more
time, could a network learn the right temporal filters for itself, straight from the
raw accelerometer trace?} This section walks the architecture end to end. It is the
``given more time'' upgrade --- a study of what a learned model buys --- and is
\emph{not} the detector that runs on the ankle.

\begin{intuition}
The freeze index is a single filter we designed by hand: ``measure energy in
$3$--$8$\,\unit{\hertz}, divide by energy in $0.5$--$3$\,\unit{\hertz}.'' A
convolutional network is a stack of \emph{many small filters whose shapes are
learned}. Instead of telling the model which frequency band matters, we show it
thousands of labelled $4$-second clips and let gradient descent sculpt the filters
that best separate freezing from walking. The freeze index is one filter we wrote;
FoGNet is sixteen-plus filters it wrote itself.
\end{intuition}

\subsection{Why a CNN, and why 1-D}

The accelerometer hands us a \emph{time series}: three channels ($a_x, a_y, a_z$)
sampled at $64$\,\unit{\hertz}. A freeze episode is a \emph{local temporal pattern} ---
a burst of rapid $3$--$8$\,\unit{\hertz} trembling --- that can appear anywhere inside
the $4$-second window. Three properties of a 1-D convolutional network match this
problem almost exactly:

\begin{itemize}
  \item \textbf{Locality.} A freeze is defined by what neighbouring samples do
  \emph{relative to each other} (a fast oscillation), not by their absolute time
  index. Convolution looks at short local stretches, which is the right unit of
  analysis.
  \item \textbf{Weight sharing.} The same small filter slides across the whole
  window, so a tremor at $0.5$\,\unit{\second} and the identical tremor at
  $3.5$\,\unit{\second} are scored by the \emph{same} weights. This is far more
  data-efficient than a dense layer that would have to relearn the pattern at every
  offset.
  \item \textbf{Translation invariance} (completed by global pooling, below): a
  freeze \emph{anywhere} in the window should count as a freeze.
\end{itemize}

We use \emph{1-D} convolution because time is the one axis we slide along; the three
accelerometer axes are treated as input \emph{channels} (like the R/G/B channels of
an image), all read together by every filter.

\begin{keyidea}
FoGNet never sees a Welch spectrum, a band, or the freeze index. Its input is the
\emph{raw} $3\times256$ window --- three axes $\times$ $4$\,\unit{\second} at
$64$\,\unit{\hertz} --- exactly the bytes the LIS3DH produced. Any notion of
``frequency band'' it ends up using is something it \emph{learned}, not something we
supplied.
\end{keyidea}

\subsection{What one convolution actually does}

Picture a short filter --- say $15$ samples wide --- as a little template laid over
the signal. You multiply the filter against the $15$ samples beneath it, add the
products into a single number, then shift one sample to the right and repeat. Slide
all the way along and you get a new, slightly shorter trace: the filter's
\emph{response}. Where the signal locally resembles the template, the response spikes;
elsewhere it stays near zero. A filter tuned to a $\sim5$\,\unit{\hertz} wobble will
light up precisely during a freeze and stay quiet during a smooth stride.

\begin{mathsbox}
For an input channel $x[n]$ and a learned filter $w$ of length $K$ (here a
\code{Conv1d} cross-correlation), the output at position $n$ is
\begin{equation}
  y[n] \;=\; b \;+\; \sum_{c=1}^{C_{\text{in}}}\;\sum_{k=0}^{K-1}
           w_{c}[k]\, x_{c}[\,n+k\,],
  \label{eq:conv}
\end{equation}
summed over all $C_{\text{in}}$ input channels. The weights $w_c[k]$ and the bias $b$
are \emph{the things learned by backprop}; nothing about them is fixed in advance.
With $C_{\text{out}}$ such filters in a layer, the output has $C_{\text{out}}$
channels --- one response trace per filter. In FoGNet's first layer $C_{\text{in}}=3$
(the axes), $K=15$, and $C_{\text{out}}=16$, so $16$ different $15$-sample templates
sweep the window simultaneously.
\end{mathsbox}

Each convolution is followed by two cheap, essential operations:
\textbf{BatchNorm}, which re-centres and re-scales each channel so training stays
numerically stable, and a \textbf{ReLU} nonlinearity, $\mathrm{ReLU}(z)=\max(0,z)$,
which clips negatives to zero. The ReLU is what lets the stack learn patterns more
intricate than the freeze index's single ratio: without it, any stack of convolutions
would collapse into one equivalent linear filter.

\subsection{Stacking blocks: wider receptive fields}

One $15$-sample filter sees only $15/64 \approx 0.23$\,\unit{\second} of signal at a
time --- enough for one or two tremor cycles, not a whole episode. Depth fixes this.
After the first two convolution blocks FoGNet applies \textbf{max-pooling} (stride
$2$), which keeps the strongest response in each adjacent pair and halves the length
of the trace. A filter in the next layer, still only a few samples wide, now spans
\emph{twice} as much real time, because each sample it reads already summarises two
original samples. Stack the effect and the \emph{receptive field} --- the slice of the
original $4$\,\unit{\second} that one deep activation can ``see'' --- grows from a
fraction of a second to a large span of the window. Early filters catch fine
texture (the wobble itself); deep filters catch coarse structure (a sustained burst
of wobble), which is exactly what distinguishes a freeze from a single misstep.

\begin{intuition}
Max-pooling is a deliberate, useful blur. By throwing away \emph{where exactly} the
peak response sat inside each little pair and keeping only \emph{how strong} it was,
the network stops caring about sub-sample timing and starts caring about whether the
freeze texture is present at all. Cheaper to compute, and more robust.
\end{intuition}

\subsection{FoGNet, end to end}

Figure~\ref{fig:fognet} draws the whole network. The data starts wide and shallow
(a long $256$-sample window, only $3$ channels) and ends short and deep (time mostly
pooled away, $64$ learned channels), before being collapsed to a single freeze
probability. Three trends run left to right: \emph{time shrinks} (pooling),
\emph{channels grow} ($3\!\to\!16\!\to\!32\!\to\!64$, more learned features), and
\emph{kernels narrow} ($15\!\to\!9\!\to\!5$, because deeper samples already cover more
real time).

\begin{figure}[H]
\centering
\resizebox{0.97\textwidth}{!}{%
\begin{tikzpicture}[font=\sffamily]
  % ---- helpers -------------------------------------------------------------
  \tikzset{
    flab/.style={font=\footnotesize\bfseries,text=navy,align=center},
    dlab/.style={font=\scriptsize,text=mute,align=center},
    op/.style={->,>=Latex,thick,steel},
  }

  % ---- tensor volumes: \cnnvol{cx}{halfw}{halfh}{front}{back} ---------------
  % INPUT  3 x 256  (tall, thin)
  \cnnvol{0}{0.22}{2.0}{fill=teal!28,draw=teal!75!black,line width=0.5pt}{fill=teal!14,draw=teal!55!black,line width=0.5pt}
  % BLOCK 1  -> 16 ch, time 256 (drawn medium tall, medium wide)
  \cnnvol{2.6}{0.55}{1.55}{fill=navy!26,draw=navy!75!black,line width=0.5pt}{fill=navy!13,draw=navy!55!black,line width=0.5pt}
  % BLOCK 2  -> 32 ch, time pooled
  \cnnvol{5.4}{0.85}{1.05}{fill=navy!26,draw=navy!75!black,line width=0.5pt}{fill=navy!13,draw=navy!55!black,line width=0.5pt}
  % BLOCK 3  -> 64 ch, time pooled more
  \cnnvol{8.1}{1.15}{0.62}{fill=navy!26,draw=navy!75!black,line width=0.5pt}{fill=navy!13,draw=navy!55!black,line width=0.5pt}
  % GAP -> 64-vector (thin tall bar)
  \cnnvol{10.4}{0.16}{1.15}{fill=purple!26,draw=purple!75!black,line width=0.5pt}{fill=purple!13,draw=purple!60!black,line width=0.5pt}
  % Dense 32
  \cnnvol{11.9}{0.13}{0.62}{fill=purple!26,draw=purple!75!black,line width=0.5pt}{fill=purple!13,draw=purple!60!black,line width=0.5pt}
  % Output 2 logits
  \cnnvol{13.2}{0.10}{0.22}{fill=freeze!26,draw=freeze!80!black,line width=0.5pt}{fill=freeze!13,draw=freeze!60!black,line width=0.5pt}

  % ---- block labels (above) -------------------------------------------------
  \node[flab] at (0,2.55) {Input\\$3\times256$};
  \node[dlab] at (0,-2.5) {raw 3-axis\\$4$\,s @ $64$\,Hz};

  \node[flab] at (2.6,2.55) {Conv block 1\\$3\!\to\!16$};
  \node[dlab] at (2.6,-2.5) {$K{=}15$, BN,\\ReLU, pool$\times2$};

  \node[flab] at (5.4,2.55) {Conv block 2\\$16\!\to\!32$};
  \node[dlab] at (5.4,-2.5) {$K{=}9$, BN,\\ReLU, pool$\times2$};

  \node[flab] at (8.1,2.55) {Conv block 3\\$32\!\to\!64$};
  \node[dlab] at (8.1,-2.5) {$K{=}5$, BN,\\ReLU};

  \node[flab] at (10.4,2.55) {GAP\\$\to 64$};
  \node[dlab,anchor=north] at (10.4,-1.45) {mean over\\time};

  \node[flab] at (11.9,2.55) {Dense\\$64\!\to\!32$};
  \node[dlab,anchor=north] at (11.9,-1.45) {ReLU,\\Drop $0.3$};

  \node[flab,text=freeze] at (13.35,2.55) {Logits\\$32\!\to\!2$};
  \node[dlab,anchor=north] at (13.35,-1.45) {softmax\\$\to P(\text{freeze})$};

  % ---- flow arrows ----------------------------------------------------------
  \draw[op] (0.45,0) -- (1.95,0);
  \draw[op] (3.25,0) -- (4.45,0);
  \draw[op] (6.35,0) -- (6.85,0);
  \draw[op] (9.35,0) -- (10.15,0);
  \draw[op] (10.6,0) -- (11.7,0);
  \draw[op] (12.06,0) -- (13.05,0);

  % ---- group brace: convolutional feature extractor ------------------------
  \draw[decorate,decoration={brace,amplitude=5pt,mirror},steel,thick]
        (-0.35,-3.15) -- (9.4,-3.15)
        node[midway,below=6pt,font=\footnotesize\bfseries,text=steel]
        {learned feature extractor (translation-aware)};
  \draw[decorate,decoration={brace,amplitude=5pt,mirror},purple,thick]
        (9.9,-3.15) -- (13.45,-3.15)
        node[midway,below=6pt,font=\footnotesize\bfseries,text=purple]
        {classifier head};
\end{tikzpicture}}
\caption{\textbf{FoGNet, a 1-D CNN ($\sim$18k parameters).} Tensor volumes shrink in
time and grow in channels left to right. Three \texttt{Conv$\to$BN$\to$ReLU} blocks
($3\!\to\!16\!\to\!32\!\to\!64$ channels, kernels $15/9/5$, max-pool after the first
two) extract learned temporal features; global average pooling collapses each of the
$64$ channels to one number, which a small dense head maps to two logits and a
softmax freeze probability. Runs in under $1$\,ms on a laptop CPU.}
\label{fig:fognet}
\end{figure}

\subsection{Global average pooling: a freeze anywhere counts}

After the third block the network holds $64$ response traces. Rather than flatten them
(which would tie every feature to a fixed time slot and explode the parameter count),
FoGNet takes the \textbf{global average} of each trace --- one mean value per channel,
a single $64$-vector.

\begin{keyidea}
Global average pooling (GAP) is what makes FoGNet \emph{translation-invariant}. A
freeze that erupts at second $1$ and the same freeze at second $3$ produce strong
responses at different \emph{positions}, but their channel \emph{averages} are nearly
identical. So a freeze ``anywhere in the window still counts'' --- which is exactly the
behaviour the windowed detector wants, and it sidesteps the dilution problem a fixed
position-sensitive layer would suffer.
\end{keyidea}

The $64$-vector then flows through the classifier head:
$\text{Linear}(64\!\to\!32)\to\text{ReLU}\to\text{Dropout}(0.3)\to\text{Linear}(32\!\to\!2)$.
Dropout randomly mutes $30\%$ of the $32$ hidden units during training, forcing the
head not to lean on any single feature --- a regulariser against overfitting the
modest dataset. The final layer emits two raw scores (\emph{logits}), and softmax
turns them into a probability:
\begin{equation}
  P(\text{freeze}) \;=\; \frac{e^{z_{1}}}{e^{z_{0}}+e^{z_{1}}},
  \label{eq:softmax}
\end{equation}
where $z_0,z_1$ are the ``no-freeze'' and ``freeze'' logits. The whole forward pass is
tiny --- about $18$k parameters --- and completes in under $1$\,\unit{\milli\second} on
a laptop CPU.

\begin{table}[H]
\centering
\caption{FoGNet layer-by-layer. Time shrinks via pooling; channel count grows; the
classifier head turns the $64$ pooled features into a freeze probability.}
\label{tab:fognet}
\small
\begin{tabular}{@{}llcc@{}}
\toprule
\textbf{Stage} & \textbf{Operation} & \textbf{Kernel} & \textbf{Channels (out)} \\
\midrule
Input        & raw 3-axis window               & ---  & $3$  \\
Conv block 1 & Conv1d, BN, ReLU, max-pool$\times2$ & $15$ & $16$ \\
Conv block 2 & Conv1d, BN, ReLU, max-pool$\times2$ & $9$  & $32$ \\
Conv block 3 & Conv1d, BN, ReLU                 & $5$  & $64$ \\
Pooling      & global average pooling          & ---  & $64$ (vector) \\
Head         & Linear, ReLU, Dropout $0.3$     & ---  & $32$ \\
Output       & Linear $\to$ softmax            & ---  & $2$ logits \\
\bottomrule
\end{tabular}
\end{table}

\subsection{How it is trained and how it scores}

FoGNet is trained on the public \textbf{Daphnet} FoG dataset (Bachlin
et~al.\ 2010). Each labelled $3\times256$ window is pushed forward to a probability;
the \emph{binary cross-entropy} loss penalises confident wrong answers; and
gradient descent / backprop nudges every filter weight and bias --- the $w_c[k]$ of
Eq.~\eqref{eq:conv} --- a little in the direction that lowers the loss. Crucially, we
evaluate with \textbf{Leave-One-Subject-Out (LOSO)} cross-validation: the network is
tested on a person it never trained on, the honest test of whether the learned filters
\emph{generalise} rather than memorise one gait.

\begin{figure}[H]
\centering
\begin{subfigure}{0.485\textwidth}
  \centering
  \includegraphics[width=\linewidth]{confusion_matrix}
  \caption{Pooled LOSO confusion ($n=8982$ windows).}
  \label{fig:cnn-cm}
\end{subfigure}\hfill
\begin{subfigure}{0.485\textwidth}
  \centering
  \includegraphics[width=\linewidth]{per_subject}
  \caption{Per-subject scores: generalisation varies by person.}
  \label{fig:cnn-persubj}
\end{subfigure}
\caption{FoGNet under Leave-One-Subject-Out. Pooling the held-out folds gives
$\mathrm{TN}=6842$, $\mathrm{FP}=1285$, $\mathrm{FN}=248$, $\mathrm{TP}=607$:
sensitivity $0.71$, specificity $0.84$, accuracy $0.83$.}
\label{fig:cnn-results}
\end{figure}

Pooling the held-out folds gives the confusion matrix in
Figure~\ref{fig:cnn-cm}: out of $8982$ windows, $\mathrm{TP}=607$,
$\mathrm{FN}=248$, $\mathrm{TN}=6842$, $\mathrm{FP}=1285$. That is a
\textbf{sensitivity of $0.71$}, a \textbf{specificity of $0.84$} and an accuracy of
$0.83$. As argued throughout this report, sensitivity and specificity --- not raw
accuracy --- are the metrics that matter when freezes are rare; the per-subject spread
in Figure~\ref{fig:cnn-persubj} is the more sobering view, showing how much
generalisation depends on the individual.

\begin{takeaway}
FoGNet \emph{learns} its temporal filters from raw accelerometer windows instead of
relying on one fixed ratio, and global average pooling lets a freeze anywhere in the
window count. It is the ``given more time'' study --- a sub-millisecond, $\sim$18k-param
network ($0.71$ sensitivity, $0.84$ specificity under LOSO) --- and it is deliberately
\emph{not} what ships on the ankle, where the transparent freeze-index detector runs.
\end{takeaway}
